\section{Conclusion and Future Work}
\label{sec:conclusion}

This study compared flat EfficientNet-B0, flat DenseNet-121, and a conditional
EfficientNet-B0 hierarchy on the same locked 3,668-image internal-test cohort.
DenseNet-121 achieved the highest accuracy, macro-F1, and weighted-F1 point
estimates, while EfficientNet-B0 retained the highest balanced accuracy.
Paired analyses supported DenseNet-121's accuracy and weighted-F1 gains
over both comparators, although its performance advantage was not uniform
across metrics.

The hierarchy did not improve upon flat EfficientNet-B0 in the reported
endpoint metrics, and the paired analysis detected no statistically
significant difference. More importantly, oracle routing increased
hierarchical macro-F1 from 0.6054 to 0.7937. Stage 1 blocked 20.079\% of malignant images and incorrectly
forwarded 22.060\% of non-malignant images, indicating that routing
errors were a substantial source of endpoint performance loss.

Neither the alternative backbone nor the hierarchy resolved rare-class
sensitivity; DenseNet-121 correctly classified only 28 of 94 SCC images.

Future work should prioritise repeated-seed evaluation, independently
validated routing improvements, larger minority-class cohorts, and external
testing on independently collected data. Until then, the systems should be
interpreted as internally evaluated experimental models rather than clinically
validated or deployment-ready tools.
